\begin{table}[t]
  \caption{Observed system performance. Deviations are engineering bounds from
    bench comparison; see Section~\ref{sec:threats} for scope.}
  \label{tab:experimental_results_table}
  \small
  \begin{tabular}{@{}lll@{}}
    \toprule
    Metric & Observation & Reference \\
    \midrule
    Temperature accuracy      & $\pm 1^{\circ}$C, $25$--$40^{\circ}$C & Digital thermometer \\
    Pulse accuracy (resting)  & $\pm 3$ BPM                           & Fingertip oximeter \\
    Button-to-display latency & $< 1.2$ s                             & Direct observation \\
    Queue ordering            & Preserved under cycling               & Enqueue/dequeue cycles \\
    Record footprint          & 73 bytes of 2~KB SRAM                 & Firmware layout \\
    Operating mode            & Offline, stand-alone                  & n/a \\
    \bottomrule
  \end{tabular}
\end{table}


We evaluated the assembled prototype along four axes, covering temperature accuracy, pulse
accuracy, control-path latency and queue robustness under sustained cycling.
Table~\ref{tab:experimental_results_table} summarises the outcomes, and
Section~\ref{sec:threats} states what they do and do not establish. All measurements were taken on
the bench prototype running the delivered offline firmware, with the ATmega32 on its internal
8~MHz oscillator and all peripherals connected as in Section~\ref{sec:hardware}. Sensing accuracy
was assessed against commercial reference instruments rather than laboratory standards, namely a
digital clinical thermometer and a fingertip pulse oximeter, with readings taken from resting
adult volunteers under normal indoor ambient conditions.

\subsection{Sensing Accuracy}
LM35 readings agreed with the reference thermometer to within $\pm 1^{\circ}$C across the range
exercised, approximately $25$--$40^{\circ}$C, which spans room ambient through febrile body
temperature. This is adequate for pre-screening and not adequate for diagnosis, since clinical
electronic thermometers are held to laboratory accuracy limits measured in tenths of a
degree~\cite{iso80601}. Three error sources plausibly account for the spread, and their remedies
differ sharply. Quantisation contributes a hard floor of $0.25^{\circ}$C per LSB, which is small
relative to the observed bound. Uncalibrated sensor tolerance contributes a fixed offset that a
one-point calibration constant in firmware would largely remove. The over-clocked ADC of
Section~\ref{sec:hardware} contributes an unquantified degradation, and of the three it is the
cheapest to fix and the one we would address first.

Pulse readings from the XD-58C agreed with the reference oximeter to within $\pm 3$~BPM under resting
conditions. The dominant source of variation was motion and transient noise during finger placement.
Readings taken immediately on contact were markedly less stable than those taken after a brief
settling period, which is why the acquisition state enforces one. This is characteristic of the
measurement principle rather than of this sensor, since photoplethysmographic amplitude and
morphology vary strongly with probe pressure, perfusion and movement~\cite{allen2007ppg}. A subtler
contributor is the time base. The sampling interval derives from delay loops clocked by the internal
RC oscillator, specified to roughly $\pm 3\%$~\cite{microchip2020}, and a systematic 3\% interval
error maps directly onto the computed rate, about 2~BPM at 70~BPM. A hardware timer with a crystal
reference would remove this term, at the cost of the crystal the design deliberately omits.

\subsection{Control-Path Latency and Queue Robustness}
Responsiveness was measured end to end, from the physical press of the doctor's button to the
appearance of the updated record on the console, so it includes debounce, queue access, string
formatting and the full I\textsuperscript{2}C refresh. Observed latency remained below 1.2~s. Queue
advancement happens once per consultation, so a sub-1.2~s response is imperceptible at that cadence.
The relevant property is not that latency is small but that it is bounded and repeatable, which
follows from the superloop architecture of Section~\ref{sec:firmware}. Note also that the debounce
lockout is a deliberate 500~ms, so roughly half the measured latency is a design constant rather than
computational cost.

Queue integrity was exercised through repeated enqueue and dequeue cycles, checking after each that
records emerged in entry order with none lost or duplicated. Ordering and record integrity held
across sustained cycling, and the two-phase button semantics behaved correctly at the boundaries
that matter in practice. A press against an empty queue produces the no-patient message rather than
a stale or undefined record, and emptying the queue mid-consultation returns the console to that
message rather than leaving the retired patient on screen. The bound here is capacity, not
correctness. With 2~KB of SRAM against a 73-byte record, achievable depth is modest and fixed at
build time, which is not the binding constraint for a waiting room holding a few tens of patients.
For a longer record history, or any requirement to survive a power interruption, the persistent
storage of Section~\ref{sec:limitations} becomes necessary rather than optional.

\subsection{Threats to Validity and Scope}
\label{sec:threats}
The figures above are easy to over-read, so we state their limits explicitly.

\textbf{The sensing results are a bench comparison, not a clinical accuracy study.} Readings were
compared against consumer-grade instruments, themselves uncalibrated against a traceable standard,
on a small group of resting adult volunteers. We therefore report observed deviation bounds rather
than statistical agreement measures such as Bland--Altman limits, which would require a
prespecified population and sample size, and the $\pm 1^{\circ}$C and $\pm 3$~BPM figures are
engineering bounds establishing fitness for triage, not validated accuracy claims. Pulse accuracy
under exertion, in patients with poor peripheral perfusion, or in paediatric and elderly subjects
was not assessed, and photoplethysmography degrades in each case~\cite{allen2007ppg}.

\textbf{No deployment study was conducted.} The claim that a guided sequential dialogue lets a
non-technical patient self-register unaided is a design argument supported by the structure of the
state machine, not a tested result. Completion rate, error rate, time-to-complete and accessibility
for patients with low literacy or impaired vision all remain open, and a field study measuring them
is the most valuable single piece of future work.

\textbf{Throughput gains are argued, not measured.} That booth-side acquisition recovers
consultation time follows from moving acquisition off the physician's critical path, and reported
consultation lengths~\cite{das2021appointment} indicate the magnitude at stake, but we did not
measure consultations with and without the system.
